\clearpage
\subsection{Grenzwerte trotz Schwingung}
\label{subsec:real_functions_squeeze}
Der Sinus besitzt für $x\to+\infty$ keinen Grenzwert:
Bei $x=\pi/2+2k\pi$ ist sein Wert $1$, bei $x=3\pi/2+2k\pi$ ist er $-1$.
Beide Arten von Eingaben werden beliebig groß. Beschränktheit allein
reicht also nicht aus.

Eine immer kleiner werdende Hülle kann die Schwingung dennoch
auf einen Grenzwert festlegen.
\begin{TheoremBox}{Einschließungsprinzip}{real_functions_squeeze}
  Gilt für die betrachtete Annäherung schließlich
  $u(x)\le f(x)\le v(x)$ und haben beide äußeren Funktionen denselben
  endlichen Grenzwert $L$, dann gilt auch
  \[
    \lim u(x)=\lim v(x)=L
    \quad\Longrightarrow\quad
    \lim f(x)=L.
  \]
  Diese Aussage heißt \emph{Einschließungsprinzip}
  \EnglishTerm{squeeze theorem}.
\end{TheoremBox}
\begin{ExampleBox}{Eine abklingende Schwingung}{real_functions_sinc_infinity}
  Bestimmen Sie das Verhalten von $f(x)=\sin x/x$ für $x\to\pm\infty$.
  \begin{ClozeArea}[34]{6}
    Aus der Beschränktheit des Sinus folgt für $x\ne0$
    \[
      \left|\frac{\sin x}{x}\right|\le\frac1{|x|}
      \quad\Longrightarrow\quad
      -\frac1{|x|}\le\frac{\sin x}{x}\le\frac1{|x|}.
    \]
    Beide Schranken gehen gegen null. Also gilt
    $\lim_{x\to\pm\infty}\sin x/x=0$.
  \end{ClozeArea}
  \begin{center}
    \CachedTikzFigure{content/resources/chapters/real_functions/tikz/functions_squeeze.tex}
  \end{center}
\end{ExampleBox}
\begin{RemarkBox}
  Hier wurde das Verhalten für große $|x|$ untersucht, nicht für $x\to0$.
  An der Stelle null hilft die Schranke $1/|x|$ nicht.
  Den dortigen Grenzwert behandeln wir bei der Differentialrechnung.
\end{RemarkBox}
